% SPDX-License-Identifier: GPL-3.0-or-later
% Copyright (C) 2026 Alessandro Greco
% 简体中文版：CJK 排版与共享技术基础设施。
\usepackage{fontspec}
\ifdefined\HCode
  % EPUB（tex4ebook 的 DVI 模式）：ctex 无法在该模式下工作，因此直接用 fontspec
  % 选择包含中文字形的 Fandol 字体，否则 DVI 中缺失的字形会从 EPUB 文本中丢失。
  % 中文名称在下方手动设置；章、部分的编号格式见 epub.tex。
  \setmainfont{FandolSong-Regular}
  \setsansfont{FandolHei-Regular}
  \setmonofont{FandolFang-Regular}
  \AtBeginDocument{%
    \renewcommand{\contentsname}{目录}%
    \renewcommand{\listfigurename}{插图}%
    \renewcommand{\listtablename}{表格}%
    \renewcommand{\figurename}{图}%
    \renewcommand{\tablename}{表}%
    \renewcommand{\appendixname}{附录}%
  }
\else
  \usepackage[UTF8,scheme=plain,heading=true,fontset=fandol]{ctex}
  \setmainfont{TeX Gyre Pagella}
  \setsansfont{TeX Gyre Heros}[Scale=MatchLowercase]
  \setmonofont{DejaVu Sans Mono}[Scale=0.85]
  \setCJKmainfont{FandolSong-Regular}
  \setCJKsansfont{FandolHei-Regular}
  \setCJKmonofont{FandolFang-Regular}
  % scheme=plain 保留标准版式，但不会翻译“Chapter”“Contents”等名称：heading=true
  % 与 \ctexset 提供中文名称和“第 N 章”“第一部分”的编号格式。
  \ctexset{
    contentsname={目录}, listfigurename={插图}, listtablename={表格},
    figurename={图}, tablename={表}, appendixname={附录},
    part={name={第,部分}, number=\chinese{part}},
    chapter={name={第,章}},
  }
  % book 类的页眉使用英文的 \@chapapp；改用 ctex 的章编号（“第1章”“附录A”）。
  \makeatletter
  \renewcommand{\chaptermark}[1]{\markboth{\ifnum\c@secnumdepth>\m@ne
    \if@mainmatter\CTEXthechapter\quad\fi\fi #1}{}}
  \makeatother
\fi
\usepackage{csquotes}
% 为 \enquote 显式指定中文引号；ctex 不提供 csquotes 中文预设。
\DeclareQuoteStyle{zhbook}{“}{”}{‘}{’}
\setquotestyle{zhbook}
\usepackage[a4paper,inner=3cm,outer=2.5cm,top=2.5cm,bottom=3cm]{geometry}
\usepackage{microtype}
\setlength{\emergencystretch}{3em}
% 中文行内几乎没有可伸缩的空白：行末遇到较长的西文或等宽字词时，默认容差下找不到
% 合适的断行点，正文便会伸出版心。放宽容差只影响原本无法排好的段落。
\tolerance=2000
\usepackage{xcolor}
\usepackage{graphicx}
\usepackage{booktabs,tabularx,longtable,multirow}
\usepackage{amsmath,amssymb}
\usepackage{siunitx}
\usepackage{enumitem}
\usepackage{tikz}
\usetikzlibrary{arrows.meta,positioning,shapes,calc,fit,backgrounds,decorations.pathreplacing,decorations.pathmorphing,matrix,chains}
\usepackage{bytefield}
\usepackage{listings}
\usepackage[most]{tcolorbox}
\usepackage[backend=biber,style=alphabetic,maxbibnames=99,date=iso,urldate=iso,seconds=false]{biblatex}
% biblatex 没有简体中文的语言模块：参考文献中的固定用语在此设为中文，日期使用 ISO 格式
% （否则会输出“In:”“visited on 09/23/2026”“Apr. 25, 2022”等英文形式）。
\DefineBibliographyStrings{english}{%
  and={和}, andothers={等}, in={载于}, url={网址}, urlseen={访问日期}, version={版本},
  volume={卷}, number={第}, pages={页}, page={页}, edition={版}, editor={编},
  editors={编}, byeditor={编}, techreport={技术报告}, phdthesis={博士学位论文}}
% 页码写作“第 303–320 页”；作者之间不用英文的“, and”；网址标签用中文。
\DeclareFieldFormat{pages}{第~#1~页}
\DefineBibliographyExtras{english}{\let\finalandcomma\empty}
\DeclareFieldFormat{url}{网址：\url{#1}}
\addbibresource{../shared/bibliography/tor.bib}
% 共享参考文献是为意大利语版编写的：这里替换其中的意大利语注释。“Consultato il ...”
% 改为 url + urldate，由 biblatex 按本版语言输出。
\DeclareSourcemap{\maps[datatype=bibtex,overwrite=true]{
  \map{
    \step[fieldsource=note, match={Consultato il 23 settembre 2026}, final]
    \step[fieldsource=howpublished, fieldtarget=url]
    \step[fieldsource=url, match=\regexp{\A\\url\{(.+)\}\z}, replace={$1}]
    \step[fieldset=urldate, fieldvalue={2026-09-23}]
    \step[fieldset=note, null]
  }
  \map{\step[fieldsource=entrykey, match=\regexp{\Atorblogwebtunnel\z}, final]\step[fieldset=note, fieldvalue={2024 年 3 月 12 日}]}
  \map{\step[fieldsource=entrykey, match=\regexp{\Aperry2018torbrowser\z}, final]\step[fieldset=note, fieldvalue={草案}]}
  \map{\step[fieldsource=entrykey, match=\regexp{\Adsa2022\z}, final]\step[fieldset=note, fieldvalue={欧盟官方公报 L 277，2022 年 10 月 27 日}]}
  \map{\step[fieldsource=entrykey, match=\regexp{\Acodiceprivacy\z}, final]\step[fieldset=note, fieldvalue={意大利个人数据保护机构（Garante per la protezione dei dati personali）发布的合并文本}]}
  \map{\step[fieldsource=entrykey, match=\regexp{\Awinter2014spoiled\z}, final]\step[fieldset=note, fieldvalue={后经其他作者参与，收录于 \emph{Privacy Enhancing Technologies Symposium}（PETS 2014）}]}
  \map{\step[fieldsource=entrykey, match=\regexp{\Alegge167\z}, final]\step[fieldset=note, fieldvalue={《意大利共和国官方公报》普通系列第 277 号，2017 年 11 月 27 日}]}
  \map{\step[fieldsource=entrykey, match=\regexp{\Anistpqc2024\z}, final]\step[fieldset=note, fieldvalue={三项后量子标准，同时发布}]}
  \map{\step[fieldsource=entrykey, match=\regexp{\Ahayden2014metadata\z}, final]\step[fieldset=note, fieldvalue={Johns Hopkins Foreign Affairs Symposium，2014 年 4 月 1 日}]}
  \map{\step[fieldsource=entrykey, match=\regexp{\Apfitzmann2010terminology\z}, final]\step[fieldset=howpublished, fieldvalue={版本 0.34}]}
}}
\DeclareBibliographyAlias{legislation}{misc}
\usepackage{imakeidx}
\makeindex[intoc,title={索引}]
\usepackage[hidelinks]{hyperref}
\usepackage[nameinlink]{cleveref}
% cleveref 未内置中文名称；为常用引用类型显式设置本地化标签。
\crefname{chapter}{章}{章}
\Crefname{chapter}{章}{章}
\crefname{section}{节}{节}
\Crefname{section}{节}{节}
\crefname{figure}{图}{图}
\Crefname{figure}{图}{图}
\crefname{table}{表}{表}
\Crefname{table}{表}{表}
\crefname{equation}{式}{式}
\Crefname{equation}{式}{式}
\crefname{appendice}{附录}{附录}
\Crefname{appendice}{附录}{附录}
% 章、节按中文习惯写作“第 9 章”“第 7.5 节”（与正文中手写的形式一致），
% 而不是逐字对应英文的“章 9”“节 7.5”；多个引用写作“第 9 章和第 13 章”。
\newcommand{\zhcrefformat}[2]{%
  \crefformat{#1}{第~##2##1##3~#2}%
  \Crefformat{#1}{第~##2##1##3~#2}%
  \crefrangeformat{#1}{第~##3##1##4--##5##2##6~#2}%
  \Crefrangeformat{#1}{第~##3##1##4--##5##2##6~#2}%
  \crefmultiformat{#1}{第~##2##1##3~#2}{和第~##2##1##3~#2}{、第~##2##1##3~#2}{和第~##2##1##3~#2}%
  \Crefmultiformat{#1}{第~##2##1##3~#2}{和第~##2##1##3~#2}{、第~##2##1##3~#2}{和第~##2##1##3~#2}%
  \crefrangemultiformat{#1}{第~##3##1##4--##5##2##6~#2}{和第~##3##1##4--##5##2##6~#2}{、第~##3##1##4--##5##2##6~#2}{和第~##3##1##4--##5##2##6~#2}%
  \Crefrangemultiformat{#1}{第~##3##1##4--##5##2##6~#2}{和第~##3##1##4--##5##2##6~#2}{、第~##3##1##4--##5##2##6~#2}{和第~##3##1##4--##5##2##6~#2}%
}
\zhcrefformat{chapter}{章}
\zhcrefformat{section}{节}
\zhcrefformat{subsection}{节}
% 其他类型（图、表、附录）的连接词：默认是英文的“and”“to”。
\newcommand{\crefpairconjunction}{和}
\newcommand{\crefmiddleconjunction}{、}
\newcommand{\creflastconjunction}{和}
\newcommand{\crefpairgroupconjunction}{和}
\newcommand{\crefmiddlegroupconjunction}{、}
\newcommand{\creflastgroupconjunction}{和}
\newcommand{\crefrangeconjunction}{--}
